% Programming with InitechBase, Appendix C: Programming Conventions of the
% Initech Accounting System -- the rules every program of the sample
% diskette follows (pubs/examples/samir/disk/*.PRG).
\chapter{Programming Conventions}
\Gates{test-pubs-build}
The programs of the Initech Accounting System are written to the same rules,
so that any programmer of the Finance and Accounting section can read and
maintain any of them.  Follow the same rules in your own programs.

\SideHead{Comment header}Every program begins with comment lines that give its
file name and title, the section that owns it, the database and index files it
expects to find open, and in a sentence or two what it does and how it is run.

\SideHead{Names}Memory variables begin with M (\texttt{MTOTAL}); parameters with P
(\texttt{PAMOUNT}).  No memory variable has the name of a field.

\SideHead{Structure}Each structure is indented by two blanks.  A program that is
more than a screen long is divided into procedures with names that say what
they do.

\SideHead{Lines}One command to a line, and no line longer than the screen.  A
long expression is built in steps, in memory variables.

\SideHead{Checking first}A program checks before it changes: that the record
exists, that it may be changed, and, for a change that cannot be undone, that
the operator confirms.

\SideHead{Questions}Character answers are taken with \cmd{ACCEPT} and converted to
capitals; amounts with \cmd{INPUT}, checked with \fn{TYPE()}; choices with
\cmd{WAIT}.

\SideHead{Dates}Dates are entered in the standard form month, day and year,
two digits each, and converted with \fn{CTOD()}.  A report that is run as of a
date takes the date from a memory variable set before the run.

\SideHead{The end}Every program and procedure ends with \cmd{RETURN}.
